\section{Document Ingestion Pipeline (DOC2FHIR)}\label{sec:doc2fhir}

A significant challenge in real-world clinical settings is the prevalence of
paper-based or scanned medical records that exist outside the FHIR ecosystem.
To bridge this gap, CFI-Care integrates a dedicated Document-to-FHIR
(DOC2FHIR) pipeline that automatically converts scanned medical documents into
structured FHIR R5 resources, making them immediately available within the
patient's record without manual data entry.

\subsection{Pipeline Overview}

The pipeline operates as a sequence of three independent microservices
orchestrated by a central gateway. A document upload from the mobile or web
client triggers the pipeline, which processes the document asynchronously and
delivers the resulting FHIR bundle to the Node.js backend's
\texttt{/api/binary} endpoint, where it is persisted as a Binary FHIR
resource linked to the relevant patient record.

\begin{enumerate}
    \item \textbf{Gateway} (FastAPI, port~8001) --- accepts uploads, manages
          the job lifecycle, and coordinates the downstream stages.
    \item \textbf{OCR Service} (PaddleOCR, port~7862) --- extracts text and
          layout metadata from the scanned PDF\@.
    \item \textbf{Mapper} (Gemma~4 via llama.cpp, port~8070) --- transforms
          normalized OCR text into a validated FHIR R5 bundle.
\end{enumerate}

\subsection{Gateway Orchestrator}

The gateway is the single entry point for document uploads. It exposes a
\texttt{POST /v1/documents/upload} endpoint that accepts a multipart payload
containing the PDF file, the target patient identifier, a document
identifier, and an upload timestamp. Upon receipt, the gateway writes the
file to a local runtime directory and creates a job record in an embedded
SQLite database, assigning the job an initial state of \texttt{QUEUED}.

An asynchronous background worker thread manages the job state machine,
advancing each job through the following states: \texttt{QUEUED} (awaiting
processing capacity), \texttt{OCR\_PROCESSING} (text extraction in
progress), \texttt{MAPPING} (FHIR bundle generation in progress),
\texttt{COMPLETED} (bundle delivered to the Node.js backend), and
\texttt{FAILED} (unrecoverable error; job written to a dead-letter directory
for inspection).

The gateway implements per-stage timeouts and a GPU semaphore to prevent
concurrent Mapper requests from exhausting inference resources. Clients may
poll job status via \texttt{GET /v1/documents/\{job\_id\}/status} and
retrieve the final bundle via \texttt{GET /v1/documents/\{job\_id\}/result}
once the job reaches the \texttt{COMPLETED} state.

\subsection{OCR Service}

The OCR service is a GPU-accelerated inference endpoint built on PaddleOCR,
supplemented by a YOLOv8 layout detection model for identifying regions of
interest within clinical documents---headers, tables, body text, and
signatures. The service accepts a PDF file and returns normalized text
alongside layout metadata that records the spatial position and
classification of each detected region. This structural context is forwarded
to the Mapper alongside the extracted text to improve FHIR field assignment
accuracy.

\subsection{Mapper (AI-Assisted FHIR Conversion)}

The Mapper service runs a medically fine-tuned Gemma~4 language model served
through \texttt{llama.cpp}, exposing an OpenAI-compatible HTTP API on
port~8070. The service receives normalized OCR text and structural metadata
and applies a domain-specific system prompt engineered to extract clinical
entities---diagnoses, medications, lab values, dates, practitioner names, and
patient identifiers---and organize them into a valid FHIR R5 Bundle.

The generated bundle is validated against the FHIR R5 schema before being
returned to the gateway. Bundles that fail schema validation are flagged as
\texttt{FAILED} and written to the dead-letter directory rather than being
delivered to the Node.js backend, preventing malformed resources from
entering the clinical record.
